% Irodov Part III §3.5 236--250, source PDF pp.135--137
% printed 126--128, original figures 3.65--3.67 reviewed.
\electromagproblem{236}{126}
有限长单层螺线管，轴向长度 $\ell$，
截面半径 $R$，单位长度匝数 $n$，
通电流 $I$。求其中心的磁感应强度。
\exsolution
轴上单个圆环距观察点 $z$ 时产生
$\mathrm dB=\mu_0 IR^2/
[2(R^2+z^2)^{3/2}]\,n\,\mathrm dz$。
中心处将各匝从
$-\ell/2$ 积分至 $+\ell/2$：
\[
B=\frac{\mu_0nIR^2}{2}
\int_{-\ell/2}^{\ell/2}
\frac{\mathrm dz}{(R^2+z^2)^{3/2}}.
\]
用积分原函数
$z/[R^2\sqrt{R^2+z^2}]$，
\[
\boxed{B=
\frac{\mu_0nI}{\sqrt{1+(2R/\ell)^2}}}.
\]
\excheck $\ell\gg R$ 时趋于
$\mu_0nI$；有限长度的场小于无限螺线管极限。

\electromagproblem{237}{126}
很长直螺线管，半径 $R$，
单位长度匝数 $n$，电流 $I$。
以端部为 $x=0$，
$x>0$ 沿轴指向管外，
$x<0$ 指向管内。
求 (a) 轴向场 $B(x)$ 与形状；
(b) 管内场离中部值
$B_0=\mu_0nI$ 相差 $\eta=1\%$
的位置至端部距离。
\exsolution
把螺线管的环匝看成
占据 $z\le0$ 的连续分布，
由毕奥—萨伐尔定律
\[
B(x)=\frac{\mu_0nIR^2}{2}
\int_{-\infty}^{0}
\frac{\mathrm dz}{[R^2+(x-z)^2]^{3/2}}.
\]
令 $u=x-z$ 后积分，得
\[
\boxed{B(x)=\frac{\mu_0nI}{2}
\left(1-\frac{x}{\sqrt{x^2+R^2}}\right)}.
\]
当 $x\to-\infty$，$B/B_0\to1$，
$x=0$ 时 $B/B_0=1/2$，
$x\to+\infty$ 时 $B\to0$。
(b) 令管内 $x=-x_0$，且
$B(-x_0)=B_0(1-\eta)$。
于是
$x_0/\sqrt{x_0^2+R^2}=1-2\eta$，
从而
\[
\boxed{x_0=
\frac{R(1-2\eta)}
{2\sqrt{\eta(1-\eta)}}
\simeq4.9R\approx5R}.
\]
\excheck 书末第237题直接给约 $5R$；
计算时必须选 $x<0$ 的\emph{管内}位置。

\electromagproblem{238}{126}
宽 $h=2.0\,\mathrm{cm}$ 的薄导电带
紧密螺旋绕成半径
$R=2.5\,\mathrm{cm}$ 的很长单层螺线管，
带内通电流 $I=5.0\,\mathrm A$。
求管内外磁场随离轴半径 $r$ 的变化。
\exsolution
薄带紧密盘绕，近似每轴向长度 $h$
对应一匝，绕向表面电流分量
$K_\varphi\simeq I/h$。
由无限长电流片磁场跳变条件，
绕向电流使管内产生
\[
\boxed{B_{\rm inside}\simeq
\frac{\mu_0 I}{h}\simeq0.31\,\mathrm{mT}},
\quad r<R.
\]
螺旋线还带一个很小的沿轴
流动的电流分量。
由每匝轴向进给 $h$、
周长 $2\pi R$，
有效轴向电流为
$I_z\simeq I h/(2\pi R)$，
故管外有方位向磁场
\[
\boxed{B_{\rm outside}(r)\simeq
\frac{\mu_0 I h}{4\pi^2Rr}},\quad r>R.
\]
更精细的几何需乘
$\cos\theta=
[1+(h/2\pi R)^2]^{-1/2}$
等微小螺旋倾角修正。
\excheck 本题的带是\emph{连续倾斜螺旋}，
不是一圈圈独立闭合圆环；
因此管外存在小的方位向磁场。

\electromagproblem{239}{126--127}
木质环形铁芯截面非常细，
均匀缠有 $N=2.5\times10^3$ 匝电流线圈，
每匝电流为 $I$。
求环形铁芯内部磁场与环形线圈
\emph{几何中心}磁场的比值。
\exsolution
设环形主半径为 $R$、截面半径很小。
环形铁芯内沿主圆周的环路，
安培定律得
$B_{\rm core}2\pi R=\mu_0NI$，
\[
B_{\rm core}\simeq\frac{\mu_0NI}{2\pi R}.
\]
环形线圈的外侧电流可用
等效中心环电流的首阶近似，
在环心产生
$B_{\rm centre}\simeq\mu_0I/(2R)$。
因此
\[
\boxed{\frac{B_{\rm core}}{B_{\rm centre}}
\simeq\frac N\pi
\simeq8.0\times10^2}.
\]
\excheck 薄截面假设允许忽略芯内径向
磁场非均匀性；木质材料磁导率近似为1。

\electromagproblem{240}{127}
长直均匀圆截面导线通电流
$I=10\,\mathrm A$。
求每单位轴向长度、贯穿导线
\emph{半个截面（从轴至表面的径向面）}
的磁通量。
\exsolution
设导线半径 $R$。
截面内的磁场
$B(r)=\mu_0Ir/(2\pi R^2)$，
方向沿方位角。
取包含导线轴的一条径向半平面，
其单位轴向长度的面积元为
$\mathrm dA=1\,\mathrm m\,\mathrm dr$，
面法线沿方位角，故
\[
\frac{\Phi}{L}=
\int_0^RB(r)\,\mathrm dr
=\frac{\mu_0I}{2\pi R^2}\frac{R^2}{2}
=\boxed{\frac{\mu_0I}{4\pi}}.
\]
代入
$\boxed{\Phi/L=1.0\,\mathrm{\mu Wb/m}}$。
\excheck 这里的“半个截面”是
\emph{沿轴伸展的径向半平面}；
若误取横截面所在的平面，
磁场切向该截面，通量必为零。

\electromagproblem{241}{127}
很长螺线管电流 $I$，
截面积 $S$，单位长度匝数 $n$。
求穿过\emph{端面平面}的磁通量。
\exsolution
远离端部的内部场
$B_0=\mu_0nI$。
端部轴向场近似为内部的一半，
$B_{\rm end}\simeq\mu_0nI/2$。
对端面面积 $S$ 积分，
在长螺线管、横截面尺度
相对螺线管长度很小时，
\[
\boxed{\Phi_{\rm end}\simeq
\frac12\mu_0nIS}.
\]
\excheck 与管内深处的截面磁通量
$\Phi_0=\mu_0nIS$ 比较，
端面约为其一半。

\electromagproblem{242}{127}
图3.65：矩形截面环形螺线管
总匝数 $N=1000$，电流
$I=1.7\,\mathrm A$，截面高
$h=5.0\,\mathrm{cm}$，
外径/内径 $\eta=1.6$。
求穿过环形芯截面的磁通量。
\exsolution
环形安培环路半径 $r$
上的磁场
$B(r)=\mu_0NI/(2\pi r)$。
截面从内半径 $a$ 到外半径 $b$，
$b/a=\eta$，
磁通量
\[
\Phi=\int_a^bB(r)h\,\mathrm dr
=\frac{\mu_0NIh}{2\pi}\ln\frac ba.
\]
因此
\[
\boxed{\Phi=
\frac{\mu_0NIh}{2\pi}\ln\eta}
=\boxed{8.0\,\mathrm{\mu Wb}}\quad\text{（约）}.
\]
\excheck 计算需要的是半径比，
外径比与内径比相同；
无需知道环形芯的绝对主半径。

\electromagproblem{243}{127}
半径为 $R=100\,\mathrm{mm}$
的细圆环，圆心磁感应强度
$B=6.0\,\mathrm{\mu T}$。
求圆环电流对应磁矩。
\exsolution
圆环圆心场
$B=\mu_0I/(2R)$，
由此求电流
$I=2RB/\mu_0$。
圆环磁矩
$p_m=I(\pi R^2)$，所以
\[
\boxed{p_m=\frac{2\pi R^3B}{\mu_0}}.
\]
代值：
\[
\boxed{p_m\simeq0.030\,\mathrm{A\,m^2}
=30\,\mathrm{mA\,m^2}}.
\]
\excheck 原书书末的纯文本提取可能把
$R^3$ 误写为 $R^2$；
由 $p_m=I\pi R^2$ 的量纲可以独立复核立方幂。

\electromagproblem{244}{127}
图3.66：$N=500$ 匝细绝缘导线
紧密绕在\emph{半个环面}上，
环面细圆截面直径
$d=5.0\,\mathrm{cm}$，
电流 $I=0.80\,\mathrm A$。
求整个绕组的磁矩。
\exsolution
每一小匝的面积约
$A=\pi(d/2)^2$，磁矩模为
$IA$，方向沿环面的切向。
半环面上的切线方向随环周
转过 $\pi$，对同一方向的投影积分：
\[
|\bm p_m|=
\frac{NIA}{\pi}
\left|\int_{-\pi/2}^{\pi/2}
\cos\theta\,\mathrm d\theta\right|
=\frac{2NIA}{\pi}.
\]
代入
$A=\pi d^2/4$，
得
\[
\boxed{p_m=\frac12NId^2
=\frac12(500)(0.8)(0.05)^2
\simeq0.50\,\mathrm{A\,m^2}}.
\]
\excheck 不能把所有小匝的磁矩当作平行，
否则会错用 $NIA$；
必须积分半环面上变化的磁矩方向。

\electromagproblem{245}{127}
图3.67：单层平面螺旋由
$N=100$ 匝紧密细线构成，
电流 $I=8\,\mathrm{mA}$，
最内、外匝半径分别
$a=50\,\mathrm{mm}$、$b=100\,\mathrm{mm}$。
求 (a) 螺旋中心磁感应强度；
(b) 总磁矩。
\exsolution
把匝数视为半径均匀连续分布，
$\mathrm dN=N\,\mathrm dr/(b-a)$。
每匝在中心贡献
$\mathrm dB=\mu_0I\,\mathrm dN/(2r)$，
故
\[
\boxed{B=
\frac{\mu_0IN}{2(b-a)}\ln\frac ba
\simeq7.0\,\mathrm{\mu T}}.
\]
磁矩为各环的
$I\pi r^2$ 叠加，
\[
p_m=I\pi\int_a^b r^2\,
\frac{N\,\mathrm dr}{b-a}
=\frac{\pi IN}{3(b-a)}(b^3-a^3).
\]
因 $b^3-a^3=(b-a)(a^2+ab+b^2)$：
\[
\boxed{p_m=
\frac{\pi IN}{3}(a^2+ab+b^2)
\simeq15\,\mathrm{mA\,m^2}}.
\]
\excheck 这是密匝连续近似；
与等间距离散100匝逐一相加只有小修正。

\electromagproblem{246}{127--128}
半径 $R$ 的薄绝缘圆盘，
一面均匀面带电密度 $\sigma$，
以角速度 $\omega$ 绕轴旋转。
求 (a) 圆心磁场；(b) 旋转电荷的磁矩。
\exsolution
半径 $r$ 宽 $\mathrm dr$
的环带带电
$\mathrm dq=2\pi\sigma r\,\mathrm dr$。
每旋转一周通过固定截面一次，
频率 $\omega/(2\pi)$，
等效环电流
$\mathrm dI=\sigma\omega r\,\mathrm dr$。
圆心处小圆环贡献
$\mathrm dB=\mu_0\mathrm dI/(2r)
=\mu_0\sigma\omega\,\mathrm dr/2$，
积分得
\[
\boxed{B(0)=\frac12\mu_0\sigma\omega R}.
\]
各小环磁矩
$\mathrm dp_m=(\pi r^2)\,\mathrm dI$，
故
\[
\boxed{p_m=\pi\sigma\omega\int_0^Rr^3\,\mathrm dr
=\frac14\pi\sigma\omega R^4}.
\]
\excheck $\sigma$ 与 $\omega$
的符号决定方向；
上述为正 $\sigma$ 的模与右手方向。

\electromagproblem{247}{128}
半径 $R=50\,\mathrm{mm}$ 的
不导电薄球壳，均匀面电荷密度
$\sigma=10.0\,\mathrm{\mu C/m^2}$，
绕过球心的轴以
$\omega=70\,\mathrm{rad/s}$ 旋转。
求球心磁感应强度。
\exsolution
用球面极角 $\theta$ 参数化，
半径 $r_\perp=R\sin\theta$
处的球面带电小环总电荷
$\mathrm dq=\sigma2\pi R^2\sin\theta\,\mathrm d\theta$。
等效小环电流
$\mathrm dI=\sigma\omega R^2\sin\theta\,\mathrm d\theta$。
该环到球心的距离都是 $R$，
球心轴向磁场分量为
\[
\mathrm dB_z=
\frac{\mu_0}{2}\frac{\mathrm dI\,(R\sin\theta)^2}
{R^3}
=\frac{\mu_0\sigma\omega R}{2}
\sin^3\theta\,\mathrm d\theta.
\]
利用
$\int_0^\pi\sin^3\theta\,\mathrm d\theta=4/3$，
\[
\boxed{B(0)=\frac23\mu_0\sigma\omega R
\simeq2.9\times10^{-11}\,\mathrm T=29\,\mathrm{pT}}.
\]
\excheck 不可把球壳当作平圆盘；
不同纬度圆环的轴向投影角不同。

\electromagproblem{248}{128}
质量 $m$、半径 $R$ 的均匀实心球
带总电荷 $q$，电荷在球内均匀分布。
绕穿过球心的轴以角速度 $\omega$ 旋转。
求磁矩 $p_m$ 及其与机械角动量 $L$ 的比。
\exsolution
均匀体电荷密度
$\rho_e=3q/(4\pi R^3)$。
绕轴距 $r_\perp$ 的体元
$\mathrm dq=\rho_e\mathrm dV$
转动产生磁矩微元
$\mathrm dp_m=\frac12\omega r_\perp^2\mathrm dq$，
因此
\[
p_m=\frac12\omega\rho_e
\int_{\rm sphere}r_\perp^2\,\mathrm dV
=\frac12\omega\rho_e\frac{8\pi R^5}{15}
=\boxed{\frac15 qR^2\omega}.
\]
均匀实心球转动惯量
$J=(2/5)mR^2$，
机械角动量 $L=J\omega$，
\[
\boxed{\frac{p_m}{L}=\frac{q}{2m}}.
\]
\excheck 磁矩可为带符号矢量，
若球带负电其磁矩方向与机械角动量相反。

\electromagproblem{249}{128}
无限长介质圆柱半径 $R$，
预先产生径向极化
$\bm P(r)=\alpha r\,\bm e_r$。
介质随后以角速度 $\omega$
绕圆柱轴旋转。
求圆柱中心磁场。
\exsolution
由极化的束缚电荷公式，
圆柱内束缚体密度
$\rho_b=-\nabla\cdot\bm P=-2\alpha$，
外侧面束缚面密度
$\sigma_b=P(R)=\alpha R$。
随材料转动，体电荷产生
$j_\varphi(r)=\rho_b\omega r=-2\alpha\omega r$，
表面电荷形成周向表面电流
$K_\varphi=\sigma_b\omega R=\alpha\omega R^2$。
无限长圆柱周向电流片外侧轴向磁场为零。
用安培边界条件，靠内表面场
$B_z(R^-)=\mu_0K_\varphi
=\mu_0\alpha\omega R^2$。
内部微分形式
$-\mathrm dB_z/\mathrm dr
=\mu_0j_\varphi=-2\mu_0\alpha\omega r$，
积分得
$B_z(r)=\mu_0\alpha\omega r^2$。
故球心（圆柱轴线）
\[
\boxed{B(0)=0}.
\]
\excheck 不能仅算带正束缚电荷的侧壁电流，
必须同时计入转动的负束缚\emph{体电荷}电流。

\electromagproblem{250}{128}
两质子以相同的平行速度
$v=300\,\mathrm{km/s}$ 运动。
求它们间磁相互作用力与静电力的比。
\exsolution
在非相对论近似下，
运动点电荷在另一质子处的磁场
$B\simeq(\mu_0/4\pi)qv/r^2$
（两速度与两质子连线垂直的情形）。
洛伦兹磁力
$F_m=qvB=\mu_0q^2v^2/(4\pi r^2)$，
静电力
$F_e=q^2/(4\pi\varepsilon_0r^2)$。
因此
\[
\boxed{\frac{F_m}{F_e}=
\mu_0\varepsilon_0v^2
=\left(\frac vc\right)^2}.
\]
代入 $v=3.0\times10^5\,\mathrm{m/s}$：
\[
\boxed{F_m/F_e\simeq1.0\times10^{-6}}.
\]
\excheck 对严格平行运动但质子相对位置
并不垂直速度时，磁力大小还包含
几何因子 $\sin^2\theta$。
原书答案采用两质子并列、连线横向的常见构型。
